%% introduction.tex
%% 

%
% $Id: $
%
%
% Compilar a .pdf con LaTeX (pdflatex)
% Es necesario instalar Beamer (paquete latex-beamer en Debian)
%

%
% Gráficos:
% Los gráficos pueden suministrarse en PNG, JPG, TIF, PDF, MPS
% Los EPS deben convertirse a PDF (usar epstopdf)
%

\documentclass{beamer}
\usetheme{Warsaw}
\usebackgroundtemplate{\includegraphics[width=\paperwidth]{format/libresoft-bg.png}}
\usepackage[spanish]{babel}
\usepackage[latin1]{inputenc}
\usepackage{graphics}
\usepackage{amssymb} % Simbolos matematicos

%\definecolor{libresoftgreen}{RGB}{162,190,43}
%\definecolor{libresoftblue}{RGB}{0,98,143}

%\setbeamercolor{titlelike}{bg=libresoftgreen}

%% Metadatos del PDF.
\hypersetup{
  pdftitle={About Libre Software},
  pdfauthor={Jesus M. Gonzalez-Barahona, Gregorio Robles},
  pdfcreator={GSyC/LibreSoft, Universidad Rey Juan Carlos},
  pdfproducer=PDFLaTeX,
  pdfsubject={Systems Course},
}
%%

\begin{document}

\title{About Libre Software}
\subtitle{Presentation of the Systems Course}
\author{Gregorio Robles}
\institute{grex@gsyc.urjc.es \\
GSyC/LibreSoft, Universidad Rey Juan Carlos}

\date{Madrid, March 25th 2011}

\frame{
\maketitle
\begin{center}
\includegraphics[width=6cm]{format/gsyc-urjc}
\end{center}
}


% Si el titulo o el autor se quieren acortar para los pies de página
% se pueden redefinir aquí:
%\title{Titulo corto}
%\author{Autores abreviado}


% LICENCIA DE REDISTRIBUCION DE LAS TRANSPAS
\frame{
~
\vspace{4cm}

\begin{flushright}
\copyright 2002-2011 Jesus M. Gonzalez-Barahona, Gregorio Robles. \\

Some rights reserved. This presentation is distributed under the
``Attribution-ShareAlike 3.0'' license, by Creative Commons,
available at \url{http://creativecommons.org/licenses/by-sa/3.0/}

\end{flushright}
}


%\section{}



%%---------------------------------------------------------------------
%%---------------------------------------------------------------------

\section{About Libre Software}

%%---------------------------------------------------------------------
\begin{frame}
\frametitle{There is a new guy in town}

\begin{itemize}
\item GNU/Linux, Apache, GNOME, KDE, LibreOffice, etc. are very
  important, but...
\item The really new thing is the libre software model:
  \begin{itemize}
  \item Unprecedented combination of collaboration and competition.
  \item Shift in emphasis from marketing to support and quality.
  \item Classical assumptions about intellectual propriety are
    questioned.
  \item End-users recover the control (instead of big software providers)
  \item A new model for a new (global, networked) world?
  \end{itemize}
\item Last years have shown the feasibility of the model.
\end{itemize}

\end{frame}


%%---------------------------------------------------------------

\begin{frame}
\frametitle{What is libre software?}

In short free software guarantees:

\begin{itemize}
\item Freedom to use
\item Freedom to study, and to adapt
\item Freedom to redistribute
\item Freedom to improve and release improvements
\end{itemize}

In other words, if you get it, you can...

\begin{itemize}
\item use it
\item study and adapt it
\item redistribute it
\item improve it and release improvements
\end{itemize}
\vspace{.2cm}
\begin{flushright}
\url{http://www.gnu.org/philosophy/free-sw.html}
\end{flushright}

\end{frame}


%%---------------------------------------------------------------

\begin{frame}
\frametitle{Free / libre / open source}

\begin{itemize}
\item The definition is from Free Software Foundation, for free
  software (``Free Software Definition'')

\item But same applies to open source software (``Open Source Definition'')

\item To avoid discussion and missinterpretations, we will ``libre software''

\item Important consequence: \\
  To be able of modifying source code, it must be available.

\item Lots of licenses: GPL, LGPL, BSD, Apache, MPL, etc.
\end{itemize}

\vspace{.2cm}

\begin{flushright}
\url{http://opensource.org/docs/osd}
\end{flushright}

\end{frame}


%%---------------------------------------------------------------

\begin{frame}
\frametitle{When libre software enters a new niche...}

\begin{itemize}
\item It can become one of the first choices (GNU/Linux in operating
  systems, Apache in WWW servers, LibreOffice in office applications, etc.)
\item It benefits from a lot of synergy (reuse of code, reuse of
  knowledge, reuse of distribution channels, etc.)
\item Users gain competitive advantage:
  \begin{itemize}
  \item Availability of source code makes improvements and
    customization possible in large scale (by in-house or
    subcontracted teams).
  \item Standardization, but maintaining competition between
    providers.
  \item No more per-use licenses.
  \item Much more and better support (ensured by competition).
  \end{itemize}
\item Competition is the name of the game.
\end{itemize}

\end{frame}


%%---------------------------------------------------------------

\begin{frame}
\frametitle{Consequences for the software industry}

The software business is changing upside down (still slowly, but
gaining momentum):

\begin{itemize} 
\item Traditional software ``manufacturers'' will have to reinvent
  themselves completely (no more per-copy incomes).
\item A whole new industry (based in support and libre development)
  will be needed as libre software gains market acceptance.
\item It allows for (and encourages) competition in support, and even
  in the evolution of a piece of software.
\item Users are benefited in several ways. Therefore, big pressure
  from end-users (including big companies) to switch to libre software.
\end{itemize}

\end{frame}

%%---------------------------------------------------------------

\begin{frame}
\frametitle{Some specific impacts}

\begin{itemize} 
\item Cost: cost model radically different from proprietary software
\item Openess: can be modified, can be inspected, can be studied
\item Distribution: new distribution channels, new methods
\item Development: ``surprising'' development models
\item Maintenance and support: true competition
\end{itemize}

Mixture of two powerful mechanisms:

\begin{itemize}
\item Competition (using the same souce base)
\item Cooperation (even non-voluntary)
\end{itemize}

\end{frame}

%%---------------------------------------------------------------

\begin{frame}
\frametitle{Different actors, different visions}

\begin{itemize}
\item End users (persons, companies, public administrations, etc.)
\item Developer (or software producer)
\item Software integration
\item Service provider
\end{itemize}

\end{frame}


%%---------------------------------------------------------------

\begin{frame}
\frametitle{Libre software for (large) end users}

\begin{itemize}
\item Libre software is not necessarily better or worse. It is 
  just different
\item In several niches, we have already excelent products and
  companies supporting them.
\item In many cases, the most cost-effective way of producing
  software.
\item Special advantages when there is interest in long-term life
  cycles, vendor independence, multiplatform support, adaption to
  evolving technologies.
\item If a powerful enough user (or group of users) needs to drive the
  technology, this is probably the best way to go.
\item Many things can be done to promote a competitive
  libre software industry in a given niche. Many benefits
  are derived of such a promotion.
\end{itemize}
\end{frame}

%%---------------------------------------------------------------

\begin{frame}
\frametitle{}

\begin{itemize}
\item 
\end{itemize}

\end{frame}


%%---------------------------------------------------------------

\begin{frame}

\begin{center}
{\LARGE
Why not learn about libre software?
}
\end{center}


\end{frame}

\end{document}
